% Chapter 06: Testing Contracts
% Acceptance tests for your Domain Kit

\chapter{Testing Contracts}
\label{ch:testing-contracts}

% ============================================================
\section{Kit Acceptance Test Structure}
\label{sec:test-structure}
% ============================================================

A Domain Kit MUST include acceptance tests that define what ``correct'' means
for the Kit. These tests verify:

\begin{enumerate}
  \item \textbf{Lowering correctness} --- The compiler produces the expected
    core YAML for Kit-authored runbooks.
  \item \textbf{Validator behavior} --- Validators correctly accept valid input
    and reject invalid input with actionable error messages.
  \item \textbf{Projection output} --- Projections produce the expected output
    for a given trace stream.
\end{enumerate}

\paragraph{Test directory structure:}

\begin{verbatim}
test/
├── fixtures/                     # Test input runbooks
│   ├── 01-basic-approval.yaml
│   ├── 02-nested-governance.yaml
│   ├── 03-evidence-capture.yaml
│   └── 04-invalid-missing-owner.yaml
└── expected/                     # Expected outputs
    ├── 01-basic-approval.lowered.yaml
    ├── 02-nested-governance.lowered.yaml
    ├── 03-evidence-capture.lowered.yaml
    └── 04-invalid-missing-owner.errors.json
\end{verbatim}

% ============================================================
\section{Test Fixture Format}
\label{sec:fixture-format}
% ============================================================

A test fixture is a Kit-authored runbook that exercises a specific feature or
validates a specific error condition.

\paragraph{Example: \texttt{01-basic-approval.yaml}}

\begin{minted}{yaml}
# Test fixture: basic approval-request lowering
apiVersion: runbook/v2+compliance/v1
meta:
  title: "Basic approval test"
  kit:
    name: gert.compliance
    version: "1.0.0"

steps:
  - id: approve-step
    type: approval-request
    description: "Simple approval with two approvers"
    approvers:
      - alice@acme.example
      - bob@acme.example
    timeout: 1h
\end{minted}

\paragraph{Expected lowered output: \texttt{01-basic-approval.lowered.yaml}}

\begin{minted}{yaml}
apiVersion: runbook/v2
meta:
  title: "Basic approval test"
  x-kit:lowered-from:
    kit: gert.compliance
    version: "1.0.0"

steps:
  - id: approve-step
    type: manual
    description: "Simple approval with two approvers"
    prompt: |
      Approval required from:
        - alice@acme.example
        - bob@acme.example
      
      Timeout: 1h
    approval:
      required: true
      approvers:
        - alice@acme.example
        - bob@acme.example
      timeout: 3600
\end{minted}

% ============================================================
\section{Running Kit Tests}
\label{sec:running-tests}
% ============================================================

The \texttt{gert kit test} command (described here as expected behavior, even
if not yet implemented) should:

\begin{enumerate}
  \item Discover all fixtures in \texttt{test/fixtures/}.
  \item For each fixture:
    \begin{enumerate}
      \item Run the Kit compiler: \texttt{cat fixture.yaml | ./bin/gert-kit-compliance compile > actual.yaml}
      \item Compare \texttt{actual.yaml} to \texttt{expected/*.lowered.yaml}.
      \item If the fixture is marked as invalid (e.g., \texttt{04-invalid-*.yaml}), assert that the compiler exits with code 1 and produces the expected error JSON.
    \end{enumerate}
  \item Report results: pass/fail for each fixture.
\end{enumerate}

\paragraph{Manual test invocation:}

Until \texttt{gert kit test} is implemented, run tests manually:

\begin{verbatim}
# Test 01: basic approval
cat test/fixtures/01-basic-approval.yaml | \
  ./bin/gert-kit-compliance compile > /tmp/actual.yaml

diff -u test/expected/01-basic-approval.lowered.yaml /tmp/actual.yaml

# Test 04: invalid missing owner (expect error)
cat test/fixtures/04-invalid-missing-owner.yaml | \
  ./bin/gert-kit-compliance compile > /tmp/errors.json
echo $?  # Should be 1 (error)

diff -u test/expected/04-invalid-missing-owner.errors.json /tmp/errors.json
\end{verbatim}

% ============================================================
\section{What to Test}
\label{sec:what-to-test}
% ============================================================

\paragraph{Lowering correctness:}

\begin{itemize}
  \item Basic step type lowering (one Kit step → one or more core steps).
  \item Field mapping (Kit fields → core fields).
  \item Governance field population (e.g., \texttt{approval-request} →
    \texttt{approval} field).
  \item Synthesized steps (if your Kit adds setup/teardown steps).
  \item Variable injection (if your Kit adds variables to \texttt{vars}).
\end{itemize}

\paragraph{Validator edge cases:}

\begin{itemize}
  \item Valid input passes validation (exit code 0, empty error array).
  \item Missing required fields produce errors.
  \item Invalid field values produce errors.
  \item Cross-step invariants are enforced (e.g., cleanup follows provision).
  \item Warnings are emitted when appropriate (non-blocking).
\end{itemize}

\paragraph{Projection output format:}

\begin{itemize}
  \item Projections produce well-formed JSON/YAML/Markdown.
  \item Projections correctly filter events (only audit-relevant events in
    audit log).
  \item Projections handle incomplete traces (if \texttt{run.completed} is
    missing).
\end{itemize}

% ============================================================
\section{Example Test Fixture for Validation}
\label{sec:validation-fixture}
% ============================================================

\paragraph{Fixture: \texttt{04-invalid-missing-owner.yaml}}

\begin{minted}{yaml}
# Test fixture: missing owner (should fail validation)
apiVersion: runbook/v2+compliance/v1
meta:
  title: "Invalid runbook: missing owner"

steps:
  - id: step-1
    type: approval-request
    description: "This should fail because meta.roles.owner is missing"
    approvers:
      - alice@acme.example
    timeout: 1h
\end{minted}

\paragraph{Expected error output: \texttt{04-invalid-missing-owner.errors.json}}

\begin{minted}{json}
[
  {
    "file": "<stdin>",
    "line": 3,
    "severity": "error",
    "message": "Compliance runbooks must declare an owner in meta.roles.owner",
    "rule": "compliance/owner-required"
  }
]
\end{minted}

\paragraph{Test invocation:}

\begin{verbatim}
cat test/fixtures/04-invalid-missing-owner.yaml | \
  ./bin/gert-kit-compliance compile 2>&1 > actual-errors.json

# Assert exit code is 1
echo $?  # Expected: 1

# Assert error JSON matches
diff -u test/expected/04-invalid-missing-owner.errors.json actual-errors.json
\end{verbatim}

% ============================================================
\section{Automating Tests with a Makefile}
\label{sec:makefile-tests}
% ============================================================

Create a \texttt{Makefile} target to run all tests:

\begin{verbatim}
.PHONY: test
test: test-lowering test-validation

.PHONY: test-lowering
test-lowering:
	@echo "Testing lowering..."
	@for fixture in test/fixtures/*.yaml; do \
		name=$$(basename $$fixture .yaml); \
		if [ -f test/expected/$$name.lowered.yaml ]; then \
			cat $$fixture | ./bin/gert-kit-compliance compile > /tmp/$$name.actual.yaml; \
			diff -u test/expected/$$name.lowered.yaml /tmp/$$name.actual.yaml || exit 1; \
			echo "✓ $$name"; \
		fi \
	done

.PHONY: test-validation
test-validation:
	@echo "Testing validation..."
	@for fixture in test/fixtures/04-invalid-*.yaml; do \
		name=$$(basename $$fixture .yaml); \
		cat $$fixture | ./bin/gert-kit-compliance compile > /tmp/$$name.actual.json 2>&1; \
		if [ $$? -eq 1 ]; then \
			diff -u test/expected/$$name.errors.json /tmp/$$name.actual.json || exit 1; \
			echo "✓ $$name"; \
		else \
			echo "✗ $$name: expected exit code 1, got 0"; \
			exit 1; \
		fi \
	done
\end{verbatim}

Run with:

\begin{verbatim}
make test
\end{verbatim}

% ============================================================
\section{Continuous Integration}
\label{sec:ci}
% ============================================================

Run Kit tests in CI (GitHub Actions, GitLab CI, etc.) to ensure regressions are
caught before merge.

\paragraph{Example GitHub Actions workflow:}

\begin{minted}{yaml}
name: Kit Tests
on: [push, pull_request]

jobs:
  test:
    runs-on: ubuntu-latest
    steps:
      - uses: actions/checkout@v3
      - name: Build Kit compiler
        run: go build -o bin/gert-kit-compliance ./cmd/compiler
      - name: Run tests
        run: make test
\end{minted}

% ============================================================
\section{Versioning Test Fixtures}
\label{sec:versioning-fixtures}
% ============================================================

As your Kit evolves, maintain backward compatibility by:

\begin{enumerate}
  \item Keeping test fixtures for all major versions (e.g.,
    \texttt{test/v1/fixtures/}, \texttt{test/v2/fixtures/}).
  \item Ensuring the compiler produces correct output for older runbooks (if
    you support multiple Kit versions).
  \item Adding new fixtures for new features without breaking old fixtures.
\end{enumerate}

This ensures that users can upgrade the Kit without rewriting their runbooks.
